\section{Related Work}
\label{sec:related}

\subsection{AI Persona Prompting for Creative Tasks}

The use of persona conditioning to shape AI creative output has a well-established
empirical literature. Park et al.\ demonstrate that generative agents endowed with
biographical persona descriptions produce more coherent long-horizon social behavior,
underscoring the functional role of identity framing in steering model behavior
\cite{park2023generative}. Dang et al.\ show that creative task framing --- including
role assignment --- affects the novelty and appropriateness of AI-generated artifacts,
though their focus is on task description rather than persona decomposition
\cite{dang2023creative}. Zhang et al.\ explore persona specification in dialogue
systems, finding that persona consistency correlates with perceived quality in
human evaluation \cite{zhang2018personas}.

The common architecture across this prior work treats personas as holistic units: a
persona is injected as a label (``you are a marketing expert'') or as a biographical
summary, and the resulting outputs are compared to a no-persona baseline. What the
literature does not do is decompose a persona into constituent cognitive traits,
measure the marginal quality contribution of each trait, and identify a hierarchical
tier structure among them. Our work is the first to perform this decomposition in a
controlled empirical setting and the first to identify orientation as a categorically
distinct and irreplaceable tier rather than one feature among many.

A separate strand of work examines system-prompt engineering for creative tasks,
including chain-of-thought framing \cite{wei2022chain}, role-playing instructions, and
few-shot exemplars \cite{brown2020language}. These methods improve output quality but
do not distinguish among the functional roles of different prompt components. Our tier
model provides a principled vocabulary for that distinction.

\subsection{Expertise Structure and Creative Cognition}

The Dreyfus model of skill acquisition identifies a progression from novice through
expert in which early stages are characterized by rule-following and domain-knowledge
application while expert performance becomes intuitive, context-sensitive, and
orientation-driven \cite{dreyfus1980mind}. The expert does not consult explicit
domain principles; they perceive the situation directly through an internalized
stance toward what matters. This phenomenological account of expertise predicts that
domain knowledge is an early-stage resource, not the defining feature of expert
output --- a prediction our empirical results confirm in the AI setting.

Chi et al.\ demonstrate that expert and novice physicists solving mechanics problems
organize the same domain knowledge differently: experts categorize by deep structure
and underlying principle, novices by surface feature and problem format
\cite{chi1981categorization}. This distinction maps directly onto our Tier 1 and Tier
3 distinction. The expert's organizational principle is an orientation (``what physical
law governs this situation?'') rather than a domain label (``this is an inclined-plane
problem''). The quality gain from orientation over domain labeling, documented here in
an AI creative setting, is the artificial analog of the expert-novice gap Chi et al.\
document in human cognition.

Csikszentmihalyi's systems model of creativity distinguishes the creator's domain
knowledge, the field's evaluative criteria, and the generative worldview that
determines what counts as a worthwhile creative direction \cite{csikszentmihalyi1997creativity}.
Our tier model maps onto this tripartite structure: Tier 3 (specificity) corresponds
to domain knowledge, Tier 2 (discipline) to the field's operational standards, and
Tier 1 (orientation) to the generative worldview. Csikszentmihalyi identifies the
worldview as the least teachable and most consequential component of creative
expertise, consistent with our finding that orientation is irreplaceable while domain
specificity is substitutable.

We note that these theoretical frameworks are post-hoc structural convergences: the tier
model was derived empirically from ablation data and only subsequently mapped to existing
theory. These mappings lend interpretive coherence but are not the model's theoretical
foundation.

\subsection{Orientation-First Design Principles}

The primacy of the receiver's experience as the organizing criterion for design has
been articulated across multiple traditions. Norman's account of human-centered design
establishes user experience as the primary constraint from which all functional
specifications are derived \cite{norman2013design}: the designer's first question is
``what will the person experience?'', not ``what should the artifact do?''. This is
the exact priority ordering we identify empirically in creative AI systems --- the
orientation question precedes and determines the evaluation of all subsequent design
decisions.

Hassenzahl's framework for experience design formalizes the distinction between
pragmatic quality (what the artifact does) and hedonic quality (how it is experienced)
and argues that hedonic quality is the primary design criterion for artifacts whose
success is defined by use rather than function alone \cite{hassenzahl2010experience}.
In puzzle design, the relevant distinction is between mechanical correctness (does the
extraction procedure work?) and experiential quality (is there a satisfying aha moment?).
Our Tier 1 orientation corresponds to hedonic primacy; our Tier 3 specificity
corresponds to pragmatic quality. The finding that removing Tier 1 causes catastrophic
failure while adding Tier 3 without Tier 1 fails to cross the quality threshold
confirms Hassenzahl's priority ordering in an empirical AI setting.

Rams's ten principles of good design converge on a single meta-principle: good design
removes the unnecessary \cite{rams1976design}. The principle is not a domain claim
--- it does not specify what is unnecessary in any particular artifact type. It is an
orientation: a stance that determines what counts as relevant information during design
decisions. Strunk and White's directive to ``omit needless words'' is the writing
analog: an orientation principle that precedes and overrides all domain-specific style
conventions \cite{strunkwhite1999}. Both examples illustrate that the most powerful
constraints in creative practice are orientation constraints, not domain constraints,
consistent with our empirical findings.

\subsection{Minimum Sufficient Conditions in Creative Systems}

The concept of the minimum sufficient set --- the smallest set of components that
produces a target outcome --- has a formal home in diagnostic reasoning. Reiter's
minimal diagnosis theory identifies the minimal set of component faults consistent
with observed system behavior as the natural unit of explanation in model-based
diagnosis \cite{reiter1987theory}. The minimum sufficient set concept from our
empirical reconstruction methodology (Tier 1 + Tier 2 alone crosses the quality
threshold; Tier 3 is enrichment, not requirement) is the creative quality analog of
minimal diagnosis.

In product development, the minimum viable product (MVP) concept captures a similar
structure: the smallest set of features that delivers the core value proposition to an
early user. The MVP is not the product minus domain features --- it is the product
minus everything except the features that constitute the core orientation toward the
user's need. Our Tier 1+2 minimum sufficient set is the MVP of creative AI prompting:
it delivers the core value without the domain enrichment that is necessary for ceiling
performance but not for threshold crossing.

Feature selection in machine learning provides a formal precedent for the distinction
between irreplaceable, load-bearing, and amplifying components. Guyon and Elisseeff's
treatment of variable selection identifies features by their marginal contribution to
predictive accuracy rather than their domain relevance, and shows that domain-relevant
features are often redundant given other features already selected \cite{guyon2003introduction}.
Our ablation methodology follows this logic: we measure quality under individual
feature removal rather than inferring contribution from domain relevance, and we find
that domain-relevant features (Tier 3) are indeed redundant given Tier 1 and Tier 2.

What no prior work provides is the identification of a tier structure in AI creative
prompts, the demonstration that this structure is symmetric across authoring and
reviewing roles, or the derivation of a minimum sufficient creative prompt formula
from empirical decomposition. These are the contributions of the present paper.
